\section*{Bilan du travail accompli}
\addcontentsline{toc}{section}{Bilan du travail accompli}

Ce projet de fin d'études a produit \nova{} : une plateforme
\emph{Software-as-a-Service} multi-locataire qui automatise la
relation client des petites et moyennes entreprises tunisiennes
via WhatsApp, et qui s'étend à une application web progressive
installable côté patient. J'ai mené ce travail selon une
méthodologie Agile Scrum disciplinée, en quatre itérations~: un
sprint préparatoire d'analyse suivi de trois sprints de
réalisation (\textbf{Sprint~1} -- Fondations, \textbf{Sprint~2} --
Surfaces opérationnelles, \textbf{Sprint~3} -- Couche
d'intelligence).

Au terme de cette démarche, j'ai conçu et implémenté~:
\begin{itemize}
  \item une colonne vertébrale Postgres multi-locataire sur
        Supabase, sécurisée par des politiques Row-Level Security
        strictes~;
  \item une intégration WhatsApp Business Cloud API avec un
        pipeline d'ingestion de messages dédupliqué~;
  \item un cerveau IA combinant Anthropic Claude Haiku
        (classification d'intention) et Claude Sonnet (OCR de
        documents manuscrits), exécutant une machine d'états de
        prise de rendez-vous multi-tours en arabe dialectal,
        français et anglais~;
  \item un tableau de bord administrateur fonctionnel avec
        calendrier, gestion des rendez-vous, des clients et des
        services~;
  \item une application web progressive patient avec onboarding
        par OTP téléphone et mode hors ligne~;
  \item une couche d'intelligence proactive (remplissage
        automatique de créneaux, vouchers d'anniversaire,
        apprentissage des durées, observation patient, avatars
        vocaux).
\end{itemize}

Toute la pile est hébergée sur Supabase et Vercel, avec un coût
d'infrastructure par locataire qui reste inférieur à trois USD
par mois en régime de croisière, confortablement en dessous du
plafond cible qui rend rentable un abonnement mensuel à 50~dinars.

\section*{Leçons méthodologiques et techniques}
\addcontentsline{toc}{section}{Leçons méthodologiques et techniques}

Travailler en Scrum avec une équipe d'une seule personne m'a
imposé un mélange inhabituel de rôles. J'ai joué Product Owner,
Scrum Master et équipe simultanément. Deux intuitions
méthodologiques méritent d'être soulignées. D'abord, la cadence
de sprint de deux à trois semaines était la bonne~: une cadence
plus courte aurait produit trop de fonctionnalités à moitié
finies, une plus longue aurait rendu les rétrospectives moins
utiles. Ensuite, écrire un fichier de migration SQL séparé par
fonctionnalité, plutôt que de modifier le schéma de base, a rendu
les \emph{rollbacks} triviaux durant les tests.

Sur le plan technique, plusieurs décisions ont payé~: JavaScript
ES~Modules vanille pour le tableau de bord (élimination d'une
étape de \emph{build} entière), fonctions Edge Supabase sur Deno
(\emph{cold-starts} sous 100~ms), prompts système JSON strict
pour Claude (sortie prévisible insérable directement en base) et
Postgres Row-Level Security comme frontière multi-locataire.

\section*{Limites du système actuel}
\addcontentsline{toc}{section}{Limites du système actuel}

\nova{} est dans un état où il peut être déployé pour une vraie
clinique demain, mais plusieurs limites doivent être honnêtement
reconnues~:
\begin{itemize}
  \item Les avatars vocaux ne sont pas actifs en production~: le
        pipeline est en place mais conservé derrière un drapeau
        de fonctionnalité en attendant un abonnement ElevenLabs
        payé.
  \item Aucune intégration de paiement en production~: une
        intégration Konnect a été prototypée mais retirée quand
        l'environnement de tests s'est révélé instable.
  \item Hébergement mono-région~: Supabase est à Francfort, le
        ping moyen depuis Tunis est d'environ 38~ms.
  \item La vérification \emph{green-tick} Meta WhatsApp Business
        n'a pas encore été obtenue.
\end{itemize}

\section*{Perspectives}
\addcontentsline{toc}{section}{Perspectives}

Plusieurs directions d'évolution sont envisagées pour la suite,
regroupées en trois catégories~: améliorations techniques,
améliorations fonctionnelles et expansion commerciale.

\subsection*{Améliorations techniques}
\begin{itemize}
  \item \textbf{Activation des avatars vocaux} une fois un
        abonnement ElevenLabs provisionné, pour servir rappels,
        confirmations et vouchers avec la voix clonée de
        l'administrateur.
  \item \textbf{Classifieur dialectal embarqué.} Entraîner un
        petit modèle dialectal tunisien capable de router 60 à
        70\,\% des intentions les plus simples vers un modèle
        gratuit, en réservant Claude pour les cas complexes.
  \item \textbf{Multi-région.} Migration vers une instance
        Supabase plus proche pour ramener le \emph{ping} depuis
        Tunis sous 20~ms.
  \item \textbf{Vérification \emph{green-tick} WhatsApp} pour
        débloquer la messagerie templatée proactive.
  \item \textbf{Normalisation des données médicales.} Remplacer le
        stockage JSONB libre des médicaments par une vraie table
        \texttt{client\_medications} avec clé étrangère vers un
        catalogue.
\end{itemize}

\subsection*{Améliorations fonctionnelles}
\begin{itemize}
  \item \textbf{Couche de paiement intégrée} (Konnect ou Stripe
        Tunisie) pour permettre aux patients de régler dans
        WhatsApp sans quitter le canal.
  \item \textbf{Téléconsultation vidéo} via WebRTC depuis la PWA,
        avec envoi automatique de l'ordonnance et facturation à
        la fin.
  \item \textbf{Tableau de bord multi-localisation} pour les
        administrateurs gérant plusieurs branches.
  \item \textbf{Réceptionniste vocal téléphonique} via Twilio
        Voice et un pipeline de transcription temps réel.
  \item \textbf{Marketplace de découverte patient} où les patients
        peuvent rechercher à travers toutes les entreprises
        participantes par service, localisation et disponibilité.
\end{itemize}

\subsection*{Expansion commerciale}
\begin{itemize}
  \item \textbf{Stratégie de communication locale.} Lancement
        public sur les réseaux sociaux et démarchage des premiers
        clients pilotes dans le tissu PME de Kairouan.
  \item \textbf{Onboarding self-service.} Réduire le temps
        d'inscription d'une entreprise à moins de dix minutes en
        s'appuyant sur les valeurs par défaut sectorielles.
  \item \textbf{Programme de pilote rémunéré.} Offrir trois mois
        gratuits aux cinq premières structures pilotes en échange
        d'études de cas chiffrées.
  \item \textbf{Expansion régionale Maghreb.} Adapter \nova{} aux
        marchés algérien et marocain dès la validation du modèle
        économique en Tunisie.
\end{itemize}

\section*{Mot de la fin}
\addcontentsline{toc}{section}{Mot de la fin}

Au-delà du commercial, ce qui me rend le plus fier dans ce
projet, c'est l'idée même~: un propriétaire d'entreprise n'a
plus besoin de payer un salaire mensuel pour un secrétaire qui
répond aux messages WhatsApp. \nova{} s'en charge 24~h/24 pour
le prix d'un repas. Pour une PME tunisienne qui se bat sur ses
marges, c'est une différence réelle, mesurable, défendable.
